\documentclass{article}
\usepackage[margin=1in]{geometry}

\begin{document}

\title{\textbf{Problem Set 2} \\ Main Tools of a Data Scientist}
\author{Jhinuk Dey}
\date{February 10, 2026}
\maketitle

\section*{Main Tools of a Data Scientist}

As discussed in class, the main tools of a data scientist include:

\begin{itemize}
    \item \textbf{Measurement}: Understanding how qualitative phenomena are quantified with data and what the data set captures
    
    \item \textbf{Statistical Programming Languages}: R, Python, and Julia for data analysis and statistical modeling
    
    \item \textbf{Web Scraping Tools}: APIs and HTML parsing packages 
    
    \item \textbf{Big Data Management Software}: Hadoop and Spark for handling Resilient Distributed Datasets (RDDs)
    
    \item \textbf{Database Languages}: SQL for data transformation, subsetting, and merging
    
    \item \textbf{Visualization Libraries}: ggplot2 (R), matplotlib (Python), Plots.jl (Julia), and Tableau for more interesting visualizations
    
    \item \textbf{Machine Learning Libraries}: mlr (R), scikit-learn (Python), Flux.jl (Julia) for predictive modeling
    
    \item \textbf{Data Wrangling Packages}: tidyverse (R) for data cleaning and manipulation
    
    \item \textbf{Version Control}: Git and GitHub for collaboration and reproducibility
    
    \item \textbf{High-Performance Computing}: Computing clusters like OSCER for handling large scale computations
    
    \item \textbf{Statistical Modeling Tools}: For testing theories, predicting behavior, and explaining causal relationships
\end{itemize}

\end{document}